%!TEX TS-program = lualatex
\documentclass[
    language=czech,
    doctype=article,
]{../../omnilatex}

\title{Vývoj vestavěných systémů s použitím jazyka Rust}
\subtitle{Bezpečnost paměti a výkon bez garbage collectoru}
\author{Ing. Martin Novák, Ph.D.}
\date{\today}
\keywords{Rust, vestavěné systémy, bezpečnost paměti, no_std, mikrokontroly}

\begin{document}
\maketitle

\begin{abstract}
Programovací jazyk Rust získává v posledních letech stále větší pozornost v oblasti vývoje vestavěných systémů. Tento článek zkoumá vhodnost jazyka Rust pro programování mikrokontrolerů a systémů s omezenými prostředky. Zaměřujeme se na vlastnosti jazyka, které jsou pro vestavěné systémy nejcennější: bezpečnost paměti bez garbage collectoru, jemnou kontrolu nad hardwarovými prostředky a vysoký výkon. Analyzujeme také současný stav ekosystému a výzvy spojené s adopcí Rustu v průmyslovém prostředí.
\end{abstract}

\section{Úvod}
Vývoj vestavěných systémů tradičně dominuje jazyk C a v menší míře C++. Tyto jazyky nabízejí nízkou úroveň abstrakce a efektivní kontrolu nad hardwarovými prostředky, avšak za cenu značného rizika chyb souvisejících s manipulací paměti. Chyby jako přetečení vyrovnávací paměti, uvolnění již uvolněné paměti nebo visící ukazatele jsou podle statistik příčinou významné části bezpečnostních zranitelností v softwaru pro vestavěné systémy. Programovací jazyk Rust přistupuje k tomuto problému zásadně jinak: jeho systém vlastnictví a půjčování zaručuje bezpečnost paměti při kompilaci, aniž by vyžadoval garbage collector.

\section{Vlastnosti jazyka Rust pro vestavěné systémy}
Systém vlastnictví a půjčování představuje fundamentální inovaci jazyka Rust. Každá hodnota v programu má přesně jednoho vlastníka, a po dobu její existence může existovat nejvýše jedna měnitelná reference nebo libovolný počet neměnitelných referencí. Kompilátor tyto pravidla vynucuje při překladu programu, čímž zabraňuje vzniku datových závodů a chyb souvisejících s neplatnou pamětí. Pro vestavěné systémy, kde korekce takových chyb za běhu programu často není možná, představuje tato garance obrovskou výhodu.

Prostředí no_std umožňuje kompilaci Rustu bez standardní knihovny, což je nezbytné pro systémy s velmi omezenými zdroji. V režimu no_std programátor ovládá alokaci paměti explicitně a může využívat pouze jádro jazyka bez závislosti na operačním systému. Kombinace no_std s atributem no_main umožňuje vytvoření binárního obrazu vhodného pro přímé spuštění na mikrokontroleru. Crate jako heapless nabízejí kolekce datových struktur pracující výhradně na zásobníku, což eliminuje potřebu dynamické alokace paměti.

Volitelná unsafe bloky poskytují nezbytný únik pro operace, které nelze bezpečně vyjádřit v rámci systému vlastnictví, jako je přímý přístup k hardwarovým registrům nebo obsluha přerušení. Klíčovým principem je izolace unsafe kódu do co nejmenších celků s jasně definovaným bezpečným rozhraním. Tím se minimalizuje plocha útoku a usnadňuje revize kódu.

\section{Ekosystém a nástroje}
Ekosystém Rustu pro vestavěné systémy se rychle rozvíjí. Projekt Embedded HAL (Hardware Abstraction Layer) definuje jednotné rozhraní pro periferie jako sériové linky, I2C sběrnice, SPI rozhraní a ADC převodníky. Knihovny implementující tyto rozhraní existují pro desítky mikrokontrolerů od výrobců jako STM, NXP, Microchip a Nordic Semiconductor. Podpora debuggerů a trasovacích nástrojů se neustále zlepšuje, přičemž probe-rs a cargo-embed umožňují pohodlný vývojový cyklus přímo z příkazové řádky.

Frameworky jako RTIC (Real-Time Interrupt-driven Concurrency) nabízejí strukturovaný přístup k obsluze přerušení a sdílení zdrojů v reálném čase s kompilací kontrolovanou stavovým automatem. Cargo-embed integruje kompilaci, nahrání firmware a sériovou konzoli do jediného nástroje, což výrazně zjednodušuje vývojový cyklus. Pro komplexnější systémy poskytuje Embassy asynchronní runtime postavený na konceptech async/await, umožňující psát neblokující kód v deklarativním stylu.

\section{Výzvy a budoucnost}
Hlavní výzvou pro širší adopci Rustu v průmyslu vestavěných systémů je nedostatek zkušených vývojářů a menší nabídka komerčních nástrojů ve srovnání s ekosystémem jazyka C. Doba kompilace Rustu je delší než u jazyka C, což může být problematické při iterativním vývoji. Některé specifické hardwarové funkce zatím nemají dobře zralou podporu v jazyce Rust.

Navzdory těmto výzvám se Rust postupně prosazuje v komerčních projektech. Společnosti jako Microsoft, Google a Amazon integrují Rust do svých vestavěných systémů a iniciativy jako Ferrocene nabízejí certifikovaný kompilátor Rustu pro bezpečnostně kritické aplikace. Budoucnost vývoje vestavěných systémů bude pravděpodobně charakterizována koexistencí C/C++ a Rustu, přičemž Rust bude postupně přebírat roli v nových projektech s vysokými nároky na bezpečnost a spolehlivost.

\section{Závěr}
Programovací jazyk Rust nabízí přesvědčivou kombinaci bezpečnosti paměti, vysokého výkonu a kontroly nad hardwarovými prostředky, která je ideální pro vývoj vestavěných systémů. Rychle rostoucí ekosystém nástrojů a knihoven, kombinovaný s jasnějšími garancemi bezpečnosti než u tradičních jazyků, činí z Rustu perspektivní volbu pro nové projekty v oblasti vestavěných systémů. Přechod na Rust vyžaduje investici do vzdělání vývojářů a adaptaci stávajících procesů, ale dlouhodobé přínosy v podobě sníženého počtu chyb a vyšší spolehlivosti tuto investici ospravedlňují.

\end{document}
